% 两层ReLU网络Rademacher复杂度证明的两步归约（定理 thm:two-layer-norm-risk）。
% 记 S(w)=sum_i sigma_i relu(w^T x_i)。第一步：sum_j a_j S(w_j) 关于 a 线性，
% 在 l1 球 sum|a_j|<=A 上的最大值落在顶点 +-A e_j，故等于 A sup_w |S(w)|。
% 第二步：relu 为 1-Lipschitz 且 relu(0)=0，收缩不等式把 E sup |S(w)| 化为 2 E sup sum sigma_i w^T x_i，
% 再由 Cauchy--Schwarz 得 2 B E||sum sigma_i x_i||_2 <= 2 B r sqrt(n)。
\begingroup
\definecolor{lt8Blue}{HTML}{7188A3}\definecolor{lt8BlueFill}{HTML}{E5EAF0}
\definecolor{lt8Green}{HTML}{829C8B}\definecolor{lt8GreenFill}{HTML}{E8EEE9}
\definecolor{lt8Yellow}{HTML}{B4A16F}\definecolor{lt8YellowFill}{HTML}{F1ECDD}
\definecolor{lt8Ink}{HTML}{3B4252}\definecolor{lt8Muted}{HTML}{6B7280}\definecolor{lt8Stroke}{HTML}{4C4D4F}
\begin{tikzpicture}[figure picture,x=1mm,y=1mm,text=lt8Ink,
  box/.style={rectangle,rounded corners=1.8pt,line width=.75pt,minimum width=44mm,minimum height=20mm,align=center,inner sep=2pt},
  lt8step/.style={-{Stealth[length=1.8mm,width=1.2mm]},draw=lt8Stroke,line width=.8pt}]
 \node[box,draw=lt8Blue,fill=lt8BlueFill] (A) at (0,0) {整个网络类\\[1mm]$\displaystyle\sup_{\sum_j|a_j|\le A,\ \|\bm w_j\|_2\le B}\ \sum_j a_j S(\bm w_j)$};
 \node[box,draw=lt8Green,fill=lt8GreenFill] (B) at (60,0) {单个隐藏单元\\[1mm]$\displaystyle A\sup_{\|\bm w\|_2\le B}\bigl|S(\bm w)\bigr|$};
 \node[box,draw=lt8Yellow,fill=lt8YellowFill] (C) at (120,0) {线性函数\\[1mm]$\displaystyle 2A\,\mathbb E\sup_{\|\bm w\|_2\le B}\sum_i\sigma_i\bm w^\top\bm x_i$};
 \draw[lt8step] (A.east)--(B.west);
 \draw[lt8step] (B.east)--(C.west);
 \node[figure note,text=lt8Muted,anchor=north,align=center] at (0,-12) {记$S(\bm w)=\sum_{i=1}^n\sigma_i\max\{0,\bm w^\top\bm x_i\}$};
 \node[figure note,text=lt8Muted,anchor=north,align=center] at (60,-12) {第一步：关于$\bm a$线性，$\ell_1$球上的\\最大值落在顶点$\pm A\bm e_j$};
 \node[figure note,text=lt8Muted,anchor=north,align=center] at (120,-12) {第二步：取期望后用ReLU的\\$1$-Lipschitz收缩去掉激活；\\$=2AB\,\mathbb E\bigl\|\sum_i\sigma_i\bm x_i\bigr\|_2\le2ABr\sqrt n$};
\end{tikzpicture}%
\endgroup%
